% Commande commune : encadré des sections encore à rédiger (à retirer de chaque section une fois écrite)
\newcommand{\aRediger}[2]{%
  \begin{tcolorbox}[colback=black!3,colframe=mycolor!60,boxrule=0.4pt,arc=1pt,left=6pt,right=6pt,top=4pt,bottom=4pt]
  \small\textbf{À rédiger par la sous-équipe #1.} #2
  \end{tcolorbox}}

\noindent Ce rapport s'inscrit dans le cadre du projet du module de Computer Vision, en troisième année du cycle
ingénieur Big Data \& Artificial Intelligence de l'École nationale des sciences appliquées de Tétouan.

\vspace{1cm}

\noindent\textcolor{mycolor}{\textbf{Coordonnées de l'école}}

\begin{itemize}
  \item Nom de l'école : École nationale des sciences appliquées de Tétouan (ENSA Tétouan)
  \item Université : Université Abdelmalek Essaâdi
  \item Adresse : BP 2222, Mhannech II, Tétouan, Maroc
  \item Site web : \url{https://ensa-tetouan.ac.ma}
\end{itemize}

\vspace{0.5cm}

\noindent\textcolor{mycolor}{\textbf{Le projet}}

\begin{itemize}
  \item Intitulé : LIFEWISE, assistant intelligent multimodal basé sur l'IoT et l'Edge AI pour une assistance proactive
  \item Encadrant : Pr. Otman AGHZOUT
  \item Responsables du projet : Karamo Sylla, Adila Yassir
  \item Équipe : 13 élèves ingénieurs répartis en six sous-équipes
  \item Dépôt du projet : \url{https://github.com/Kaaramo/lifewise}
\end{itemize}

\vspace{0.5cm}

\noindent\textcolor{mycolor}{\textbf{Période}}

\begin{itemize}
  \item Année universitaire 2026-2027, de novembre 2026 à la date de rendu fixée par l'encadrant
\end{itemize}
